\section{Introdução}
O projeto Night Heist combina um documento de design elaborado por estudantes com um jogo implementado por agente ChatGPT, segundo a declaração do responsável pelo projeto; o material disponível inclui especificação, código e documentação de desenvolvimento, permitindo examinar a relação entre esses artefatos \fonte. Essa configuração desloca o interesse da simples existência de um jogo funcional para a identificação de quais decisões estavam documentadas e de quais detalhes foram estabelecidos durante a implementação \fonte.

A literatura recente situa o uso de IA generativa em tarefas de engenharia de software, incluindo priorização de requisitos em contexto educacional, e também problematiza a relação entre assistência automatizada e competência de programadores iniciantes (Queiroz; Lima, 2025; Prather et al., 2024). Esses antecedentes justificam examinar artefatos e processos de especificação sem assumir que a produção final, isoladamente, demonstre aprendizagem ou domínio da programação (Prather et al., 2024).

O problema concreto deste caso consiste na distância entre enunciados do GDD e parâmetros executáveis: o documento descreve retorno ao ponto inicial, temporizador e terminal de hack, enquanto o código precisa definir condições espaciais, valores temporais e efeitos dessas interações \fonte. Sem rastreabilidade, uma decisão adicionada para tornar a regra executável pode ser confundida com requisito original, e uma divergência pode ser interpretada como defeito sem consulta ao contexto de design \fonte.

A pergunta de pesquisa é: \emph{como as decisões explicitadas no GDD discente se relacionam com a implementação do Night Heist e quais correspondências, divergências e lacunas de especificação podem ser documentadas por inspeção estática?} A formulação delimita o objeto à relação entre artefatos, conforme o protocolo do projeto, e não à avaliação individual dos estudantes \rep.

O objetivo geral é analisar essa relação mediante uma matriz de rastreabilidade com critérios observáveis, localização das evidências e distinção entre estados de verificação \rep. Os objetivos específicos são extrair unidades do GDD, relacioná-las ao código, classificar sua correspondência, identificar parâmetros não especificados e registrar os limites de atribuição das decisões \rep. Neste manuscrito, esses objetivos são operacionalizados por um piloto de seis unidades, sem apresentar o piloto como inventário completo do documento \rep.

A contribuição apresentada é um procedimento documental explicitado e acompanhado de exemplos auditáveis, que pode orientar a revisão humana do projeto e a continuidade da análise \rep. Sua relevância científica reside em tornar examinável a tradução de especificação humana para implementação assistida por IA; a originalidade em relação a outros procedimentos e a utilidade em outros projetos ainda precisam ser investigadas, sem alegação de método validado ou superioridade da IA (Queiroz; Lima, 2025; Night Heist, 2026).
